\section{导读：转折之中的 139 岁巨人}

本节先定位这期。奔驰是燃油车时代最具象征性的公司之一，康林松在 2019 年成为全球 CEO，恰好站在电动化、智能化、中国新能源竞争和豪华车重新定义的交汇点上。张小珺把他称作“转型期 CEO”，因为这不是一次普通产品换代，而是奔驰 139 年历史上最重大的技术和组织变革之一。

\begin{importantbox}{本期核心命题}
奔驰的挑战不是在“豪华”和“电动化”之间二选一，而是在中国速度、全球研发、AI 汽车、MBOS、品牌保值和财务可持续之间重新组合。康林松的回答始终围绕一个身份：他不是只追短期销量的经理人，而是 Mercedes-Benz 品牌的监护人。
\end{importantbox}

\begin{figure}[H]
\centering
\includegraphics[width=0.92\textwidth]{mercedes-transition-map.png}
\caption{139 岁奔驰的转型地图：中国速度、电动化、AI 和豪华品牌同时重构。自制概念图，依据 00:00:12--00:03:42 与全片内容整理。}
\end{figure}

\begin{knowledgebox}{读图：奔驰转型不是单变量问题}
图中 Mercedes 位于中心，周围是 China、EV、AI、Luxury、CEO 和 R\&D。读这张图时要避免把问题简化为“奔驰电动车卖得好不好”。真正问题是：一家百年豪华车公司如何在多重变量同时变化时，仍然保住品牌、技术和经营能力。
\end{knowledgebox}

\subsection{阅读路线}

本节给出读法。EP100 不是纯 AI 访谈，但它属于广义互联网与科技转型队列，因为汽车正在变成软件、算力、AI、智能座舱和全球研发网络共同定义的产品。正文将按四条线展开：中国市场为什么成为全球转型压力测试；豪华战略和电动化为什么不是对立；MBOS 与 AI 如何让奔驰重新定义技术控制权；转型期 CEO 如何在职业经理人结构里保持创业精神。

\noindent\begin{tabularx}{\textwidth}{p{0.18\textwidth}p{0.34\textwidth}X}
\toprule
阅读问题 & 访谈中的材料 & 要形成的判断 \\
\midrule
中国市场为何关键 & 中国速度、智能座舱、腾讯合作、140 亿人民币投资 & 中国不只是销售市场，也是技术和研发源头。 \\
豪华和电动化是否冲突 & 豪华车战略、纯电 G、CLA、保值和现金流 & 豪华提供利润和品牌，电动化是全组合转型方向。 \\
AI 汽车如何成立 & MBOS、车载算力、虚拟助手、DeepSeek & 汽车智能化需要架构控制和伙伴集成。 \\
CEO 如何转型巨头 & 51/49 决策、品牌监护人、创业精神、精简流程 & 转型管理不是押一个指标，而是持续校准组织。 \\
\bottomrule
\end{tabularx}

\begin{knowledgebox}{课堂提示：汽车正在变成软件定义的工业系统}
这期与 Agent 创业访谈不同，但底层问题相似：谁掌控环境，谁整合技术，谁定义体验。奔驰谈 MBOS，和 EP101 谈 YouWare 的容器/环境，都是在讨论“智能能力如何被放进一个可控系统”。
\end{knowledgebox}

\subsection{本章小结}

EP100 的主线是：老牌工业公司面对 AI 和电动化时，不能只讲历史，也不能只追新技术。它必须把品牌、技术、研发、财务和组织速度同时重构。

